\documentclass[runningheads]{llncs}

\usepackage[T1]{fontenc}
\usepackage{graphicx}
\usepackage{amsmath}
\usepackage{amssymb}
\usepackage{algorithm}
\usepackage{algpseudocode}
\usepackage{booktabs}
\usepackage{multirow}
\usepackage{listings}
\usepackage{xcolor}
\usepackage{hyperref}
\usepackage{tikz}
\usetikzlibrary{shapes.geometric, arrows, positioning, fit, calc}

\lstset{
  basicstyle=\ttfamily\footnotesize,
  breaklines=true,
  frame=single,
  keywordstyle=\color{blue},
  commentstyle=\color{gray},
  stringstyle=\color{red!60!black},
}

\begin{document}

\title{UnderCurrent: A Dual-Track Autonomic Control Plane for Risk-Aware Container Remediation}

\titlerunning{UnderCurrent: Dual-Track Autonomic Container Remediation}

\author{Pranav Negi\inst{1} \and
Reema\inst{1}}

\authorrunning{P. Negi and Reema}

\institute{Department of Computer Science and Engineering\\
\email{\{pranav.negi, reema\}@university.edu}}

\maketitle

%------------------------------------------------------------------
\begin{abstract}
Modern containerized workloads exhibit fundamentally distinct failure
signatures depending on whether they are stateless or stateful.
Stateless containers typically manifest failures through process crashes
and TCP connectivity loss, whereas stateful containers degrade through
block I/O latency accumulation, VFS write errors, and volume mount
disruptions. Existing autonomic systems apply uniform remediation
policies to both workload classes, incurring either under-reaction to
stateful I/O degradation or over-reaction to transient stateless
network events.

We present \textbf{UnderCurrent}, a dual-track autonomic control plane
that operates two parallel decision engines---Type~S for stateless
containers and Type~F for stateful containers---each governed by
workload-appropriate confidence scoring and remediation logic. Both
tracks execute every reconcile cycle on simultaneous real and shadow
paths, enabling continuous divergence measurement between executed
decisions and counterfactual alternatives. Type~F introduces a formal
five-state finite state machine (FSM) to gate irreversible actions,
preventing checkpoint-and-restart operations until a container has
progressed through a mandatory audit phase. A secondary WASM policy
sandbox provides an independent enforcement layer, blocking actions
that violate FSM-sequence constraints even if the primary DAG gate is
bypassed. We define six quantitative research metrics---misclassification
rate, divergence rate, mean time to recovery, reversibility ratio,
explanation completeness, and unsafe-action suppression rate---and
demonstrate their computability from append-only JSONL audit logs.
An evaluation framework combining cross-controller scenario replay and
threshold sensitivity analysis validates that UnderCurrent achieves
near-zero misclassification and high reversibility ratios across the
full fault severity spectrum.

\keywords{Autonomic computing \and Container orchestration \and
Finite state machines \and Kernel observability \and eBPF \and
Dual-track control plane \and Shadow execution \and
Risk-aware remediation}
\end{abstract}

%==================================================================
\section{Introduction}
\label{sec:intro}

Containerized microservice deployments have become the dominant
infrastructure paradigm for cloud-native applications
\cite{burns2016borg,merkel2014docker}. Kubernetes and similar
orchestrators provide declarative workload management but delegate
\emph{reactive} fault handling---the decision of \emph{how} to repair
a degrading container in real time---to operators or to coarsely tuned
liveness probes. Such probes treat all container failure as identical:
a threshold crossing triggers a restart regardless of whether the
container holds ephemeral web traffic or a replicated database journal.

This conflation is problematic.  A stateless HTTP worker that crashes
should be restarted immediately; the same action applied to a stateful
Postgres instance mid-transaction risks data corruption and journal
inconsistency. Conversely, suppressing action on a network-partitioned
stateful node while its I/O queue grows silently can compound the
eventual recovery cost. The root asymmetry is that stateless and
stateful workloads emit \emph{categorically different} kernel-level
failure signals, and no single remediation policy maps well to both
signal spaces.

We address this gap with \textbf{UnderCurrent}, a dual-track autonomic
control plane. Our key contributions are:

\begin{enumerate}
  \item \textbf{Workload-typed dual control plane.} Two independent
    tracks---Type~S (stateless) and Type~F (stateful)---each instrument
    the kernel via eBPF probes tuned to their respective failure
    signatures and execute remediation actions calibrated to the risk
    profile of their workload class.

  \item \textbf{Formal FSM-gated remediation for stateful workloads.}
    Type~F embeds a five-state FSM
    (\textsc{Healthy}$\to$\textsc{Degraded}$\to$\textsc{Audited}
    $\to$\textsc{Repairing}$\to$\textsc{Recovered}$\to$\textsc{Healthy})
    that enforces mandatory audit before checkpoint-and-restart, preventing
    the most disruptive action from being executed prematurely.

  \item \textbf{Real/shadow dual-path execution.} Every reconcile cycle
    executes both a live decision path and a shadow path on a throwaway
    FSM instance, producing a continuous divergence signal that quantifies
    how much the real controller's accumulated state history would have
    changed its decision relative to a stateless alternative.

  \item \textbf{WASM policy sandbox.} An optional WebAssembly policy
    module provides a second, independent gate that re-enforces
    FSM-sequence constraints at the action-execution boundary,
    independent of the reconcile-layer DAG.

  \item \textbf{Six formal research metrics.} We define misclassification
    rate, divergence rate, MTTR, reversibility ratio, explanation
    completeness, and unsafe-action suppression rate---all computable
    from append-only JSONL trace logs---and provide a cross-controller
    evaluation framework and threshold sensitivity analyzer.
\end{enumerate}

The remainder of this paper is organized as follows.
Section~\ref{sec:background} reviews related work.
Section~\ref{sec:architecture} presents the system architecture.
Sections~\ref{sec:typeS} and~\ref{sec:typeF} describe the stateless
and stateful tracks in detail.  Section~\ref{sec:metrics} defines the
research metrics and evaluation framework.
Section~\ref{sec:evaluation} reports experimental results.
Section~\ref{sec:discussion} discusses limitations and future work.
Section~\ref{sec:conclusion} concludes.

%==================================================================
\section{Background and Related Work}
\label{sec:background}

\subsection{Autonomic Computing}

The autonomic computing vision of IBM's MAPE-K (Monitor, Analyze, Plan,
Execute over a shared Knowledge base) loop \cite{kephart2003vision}
established the canonical framework for self-managing systems.
UnderCurrent maps directly onto MAPE-K: eBPF event listeners perform
monitoring, confidence scoring performs analysis, the DAG engine
performs planning, and the actions module performs execution, with
the shared constants file and trace schema serving as the knowledge base.

\subsection{Kubernetes Autoscaling and Remediation}

Kubernetes Horizontal Pod Autoscaler (HPA) and Vertical Pod Autoscaler
(VPA) address capacity management but not failure-mode-specific
remediation \cite{burns2016borg}. Custom controllers implemented via the
Operator pattern \cite{dobies2020kubernetes} allow workload-specific
logic but require per-application development effort.
UnderCurrent generalizes workload-type classification into two broad
tracks (stateless/stateful) without application-specific knowledge.

\subsection{eBPF-Based Observability}

BCC \cite{bcc2016}, bpftrace \cite{bpftrace2019}, and Cilium
\cite{cilium2018} demonstrate that eBPF can provide low-overhead
kernel-level instrumentation for network, I/O, and process events.
UnderCurrent uses eBPF tracepoints and kprobes to capture
\texttt{sched\_process\_exit}, \texttt{tcp\_connect} failures,
block I/O completions, and VFS errors without modifying container
images or introducing sidecar overhead.

\subsection{Formal Methods in Systems}

TLA+ and Alloy have been applied to verify distributed protocols
\cite{lamport2002specifying}, but runtime enforcement of formally
specified state machines in production systems remains limited.
Closest to our approach is Stateright \cite{stateright2021}, a
model-checking framework for distributed systems, and work on
runtime verification \cite{leucker2009brief}. UnderCurrent's FSM
operates as a lightweight runtime enforcer rather than an offline
verifier.

\subsection{Shadow Traffic and Canary Testing}

Dark launches and shadow traffic \cite{fowler2014dark} are widely used
in web service deployments to validate new code paths. UnderCurrent
applies the same principle at the remediation-decision level: shadow
executions compute but do not apply decisions, producing continuous
counterfactual divergence measurements at zero operational risk.

\subsection{WebAssembly Policy Enforcement}

OPA (Open Policy Agent) \cite{opa2016} and WASM-based policy enforcement
in service meshes \cite{Envoy2022wasm} demonstrate the viability of
sandboxed policy modules as independent safety gates.
UnderCurrent's WASM sandbox applies this pattern to remediation
action validation, using \texttt{wasmtime} as the runtime with a
Python fallback for portability.

%==================================================================
\section{System Architecture}
\label{sec:architecture}

\subsection{Overview}

Figure~\ref{fig:architecture} shows the top-level architecture of
UnderCurrent. Two independent control plane tracks share a common
threshold configuration ({\small\texttt{shared/constants.py}}),
a versioned trace schema ({\small\texttt{shared/trace\_schema.py}}),
and a unified evaluation layer.

\begin{figure}[t]
\centering
\begin{tikzpicture}[
  node distance=0.6cm and 0.8cm,
  box/.style={rectangle, draw, rounded corners=3pt,
              minimum width=2.0cm, minimum height=0.7cm,
              font=\footnotesize, align=center},
  arrow/.style={->, >=stealth, thick},
  track/.style={rectangle, draw=gray!60, dashed, rounded corners=5pt,
                inner sep=6pt}
]

%% Type S row
\node[box, fill=blue!15]  (evS)  {eBPF\\Events (S)};
\node[box, fill=blue!15, right=of evS]  (stS)  {State\\Store (S)};
\node[box, fill=blue!15, right=of stS]  (csS)  {Confidence\\Score (S)};
\node[box, fill=blue!15, right=of csS]  (rcS)  {DAG\\Reconcile};
\node[box, fill=blue!15, right=of rcS]  (acS)  {Actions\\(S)};
\node[box, fill=blue!15, right=of acS]  (lgS)  {Trace\\Logger};

%% Type F row
\node[box, fill=orange!20, below=1.4cm of evS] (evF)  {eBPF\\Events (F)};
\node[box, fill=orange!20, right=of evF] (stF)  {State\\Store (F)};
\node[box, fill=orange!20, right=of stF] (csF)  {Confidence\\Score (F)};
\node[box, fill=orange!20, right=of csF] (rcF)  {FSM-gated\\Reconcile};
\node[box, fill=orange!20, right=of rcF] (acF)  {Actions\\+ WASM};
\node[box, fill=orange!20, right=of acF] (lgF)  {Trace\\Logger};

%% FSM box
\node[box, fill=green!15, above=0.5cm of rcF] (fsm) {Container\\FSM};

%% Shared
\node[box, fill=gray!15, below=0.7cm of rcF] (sh)  {Shared Constants\\(thresholds, schema)};

%% Arrows Type S
\draw[arrow] (evS) -- (stS);
\draw[arrow] (stS) -- (csS);
\draw[arrow] (csS) -- (rcS);
\draw[arrow] (rcS) -- (acS);
\draw[arrow] (acS) -- (lgS);

%% Arrows Type F
\draw[arrow] (evF) -- (stF);
\draw[arrow] (stF) -- (csF);
\draw[arrow] (csF) -- (rcF);
\draw[arrow] (rcF) -- (acF);
\draw[arrow] (acF) -- (lgF);
\draw[arrow] (fsm) -- (rcF);
\draw[arrow] (rcF) -- (fsm);

%% Track boxes
\begin{scope}[on background layer]
  \node[track, fit=(evS)(lgS), label=above:{\footnotesize\textbf{Type S (Stateless)}}] {};
  \node[track, fit=(evF)(lgF)(fsm), label=below:{\footnotesize\textbf{Type F (Stateful)}}] {};
\end{scope}

\end{tikzpicture}
\caption{UnderCurrent dual-track control plane architecture.
Type~S (blue) and Type~F (orange) operate independently on
workload-appropriate eBPF probes. Type~F additionally maintains a
formal FSM that gates remediation actions. Both tracks share
threshold constants and a common trace schema.}
\label{fig:architecture}
\end{figure}

\subsection{Shared Configuration and Schema}

All decision thresholds reside in a single shared module
(Table~\ref{tab:thresholds}).  Schema versioning follows semantic
versioning: major increments signal backward-incompatible field
removals, minor increments signal new field additions, and patch
increments signal documentation changes.  The current schema version is
1.2.0. Stateless consumers reading stateful trace files silently ignore
unknown fields, ensuring backward compatibility.

\begin{table}[t]
\centering
\caption{Shared decision thresholds (\texttt{shared/constants.py}).}
\label{tab:thresholds}
\begin{tabular}{llll}
\toprule
\textbf{Constant} & \textbf{Value} & \textbf{Applicable Track} & \textbf{Triggers} \\
\midrule
\texttt{RESTART\_THRESHOLD}    & 0.80 & Type S        & restart container \\
\texttt{RESCHEDULE\_THRESHOLD} & 0.95 & Type S        & reschedule to new node \\
\texttt{FLUSH\_THRESHOLD}      & 0.50 & Type F        & flush I/O queue \\
\texttt{REPAIR\_THRESHOLD}     & 0.80 & Type F        & checkpoint-and-restart \\
\texttt{ESCALATE\_THRESHOLD}   & 0.95 & Both          & escalate / reschedule \\
\texttt{DEGRADE\_THRESHOLD}    & 0.50 & Type F (FSM)  & FSM degrade entry \\
\texttt{WINDOW\_SECONDS}       & 60   & Both          & event sliding window \\
\bottomrule
\end{tabular}
\end{table}

\subsection{Dual-Path (Real/Shadow) Execution}

On every reconcile cycle both tracks execute \emph{two} complete
decision pipelines.  The \emph{real path} uses the live state store
and, for Type~F, the live FSM (mutations persist across cycles).  The
\emph{shadow path} uses the same state store but a throwaway FSM
instance, so shadow decisions are never contaminated by prior real-path
history. Shadow actions are computed but never executed. Both paths
produce trace entries that are appended to the same JSONL log with a
\texttt{"mode"} field of \texttt{"real"} or \texttt{"shadow"}.  Per-container
divergence (real action $\neq$ shadow action) is recorded in a
separate \texttt{divergence\_log.jsonl}.

%==================================================================
\section{Type S: Stateless Control Plane}
\label{sec:typeS}

\subsection{Kernel Probes and Signal Space}

Type~S captures two kernel signals via eBPF:
\begin{itemize}
  \item \textbf{process\_exit} --- instrumented via the
    \texttt{sched:sched\_process\_exit} tracepoint.  Indicates an
    unplanned container process termination.
  \item \textbf{tcp\_connect\_fail} --- instrumented via a kprobe on
    \texttt{tcp\_connect}. Indicates sustained inability to establish
    outbound TCP connections, characteristic of network partition or
    upstream service failure.
\end{itemize}

Events are buffered in a 60-second sliding window per container.
Window maintenance prunes events older than $t_\text{now} - 60$ on
every read, ensuring that confidence scores reflect only recent history.

\subsection{Confidence Scoring}

Type~S computes two independent signal scores and combines them via
the maximum operator:

\begin{equation}
  c_S(\text{container}) = \max\bigl(s_\text{exit},\ s_\text{tcp}\bigr)
  \label{eq:score_s}
\end{equation}

\noindent\textbf{Process exit score.}  Any process\_exit event in the
window is sufficient for high confidence:
\begin{equation}
  s_\text{exit} =
  \begin{cases}
    0.95 & \text{if } |\{e \in W : e.\text{event} = \texttt{process\_exit}\}| > 0 \\
    0.0  & \text{otherwise}
  \end{cases}
  \label{eq:score_exit}
\end{equation}

\noindent\textbf{TCP connect-fail score.}  Frequency-scaled, saturating at
five events:
\begin{equation}
  s_\text{tcp} = 0.90 \cdot \min\!\left(\frac{|\{e \in W :
    e.\text{event}=\texttt{tcp\_connect\_fail}\}|}{5},\ 1.0\right)
  \label{eq:score_tcp}
\end{equation}

The maximum combination ensures that a single process exit
(catastrophic signal) dominates any number of TCP failures
(degradation signal).

\subsection{Decision DAG}

The Type~S reconciler implements a three-tier threshold DAG:

\begin{equation}
  \text{action}_S =
  \begin{cases}
    \texttt{reschedule} & c_S \geq 0.95 \\
    \texttt{restart}    & 0.80 \leq c_S < 0.95 \\
    \texttt{no\_action} & c_S < 0.80
  \end{cases}
  \label{eq:dag_s}
\end{equation}

Each decision is annotated with a \texttt{dag\_pattern} label
(\texttt{nominal}, \texttt{network\_degraded}, \texttt{process\_failure})
and the list of kernel signals that contributed to the score.

%==================================================================
\section{Type F: Stateful Control Plane}
\label{sec:typeF}

\subsection{Kernel Probes and Signal Space}

Type~F captures five kernel signals:

\begin{itemize}
  \item \textbf{blk\_io\_latency} --- block I/O completion with a
    \texttt{latency\_us} field.  High-latency events ($\geq 10{,}000~\mu$s)
    indicate I/O subsystem stress.
  \item \textbf{vfs\_write\_error} / \textbf{vfs\_read\_error} ---
    VFS-layer write and read syscall failures.  Indicate filesystem
    corruption or storage backend unavailability.
  \item \textbf{volume\_mount\_lost} / \textbf{volume\_mount\_restored} ---
    Volume mount state transitions.  A lost mount is the highest-severity
    Type~F signal.
\end{itemize}

The synthetic event generator produces these signals with weights
$[0.55, 0.18, 0.12, 0.05, 0.10]$ respectively, with blk\_io\_latency
events using a bimodal latency distribution (80\% below threshold,
20\% above) to model realistic I/O jitter.

\subsection{Confidence Scoring}

Three independent scorers are combined via MAX:

\begin{equation}
  c_F = \max\bigl(s_\text{io},\ s_\text{vfs},\ s_\text{vol}\bigr)
  \label{eq:score_f}
\end{equation}

\noindent\textbf{I/O latency score.}  Scaled by count of high-latency events:
\begin{equation}
  s_\text{io} = 0.85 \cdot \min\!\left(
    \frac{|\{e \in W : e.\text{event}=\texttt{blk\_io\_latency},
      e.\text{lat} \geq 10^4\}|}{5},\ 1.0\right)
  \label{eq:score_io}
\end{equation}

\noindent\textbf{VFS error score.}  Scaled by combined read+write error count:
\begin{equation}
  s_\text{vfs} = 0.88 \cdot \min\!\left(
    \frac{|\{e \in W : e.\text{event} \in
      \{\texttt{vfs\_write\_error}, \texttt{vfs\_read\_error}\}\}|}{3},
    \ 1.0\right)
  \label{eq:score_vfs}
\end{equation}

\noindent\textbf{Volume mount loss score.}  State-based; determined by
temporal ordering of mount events:
\begin{equation}
  s_\text{vol} =
  \begin{cases}
    0.97 & \text{if last mount event in } W \text{ is \texttt{volume\_mount\_lost}} \\
    0.0  & \text{otherwise (mounted or no event)}
  \end{cases}
  \label{eq:score_vol}
\end{equation}

Note that $s_\text{vol} = 0.97$ immediately triggers escalation
($\geq 0.95$), reflecting the high criticality of a missing data volume.

\subsection{Finite State Machine}

The Type~F FSM enforces a mandatory degradation sequence before
checkpoint-and-restart is permitted.  Formal transition function
$\delta : \Sigma \times T \to \Sigma$ where $\Sigma$ is the state set
and $T$ is the trigger set is defined in Table~\ref{tab:fsm}.

\begin{table}[t]
\centering
\caption{ContainerFSM formal transition table.}
\label{tab:fsm}
\begin{tabular}{lll}
\toprule
\textbf{From State} & \textbf{Trigger} & \textbf{To State} \\
\midrule
\textsc{Healthy}   & degrade         & \textsc{Degraded} \\
\textsc{Degraded}  & audit           & \textsc{Audited}  \\
\textsc{Degraded}  & degrade         & \textsc{Degraded} (reconfirm) \\
\textsc{Audited}   & approve\_repair  & \textsc{Repairing} \\
\textsc{Audited}   & degrade         & \textsc{Degraded} \\
\textsc{Repairing} & recover         & \textsc{Recovered} \\
\textsc{Repairing} & degrade         & \textsc{Degraded} \\
\textsc{Recovered} & heal            & \textsc{Healthy} \\
\textsc{Recovered} & degrade         & \textsc{Degraded} \\
\midrule
\multicolumn{3}{l}{\textit{All other (state, trigger) pairs are invalid (rejected).}} \\
\bottomrule
\end{tabular}
\end{table}

New containers default to \textsc{Healthy}.  Invalid triggers are
silently rejected, preserving the current state.  All transitions are
recorded in a per-container history log for audit purposes.

\subsection{FSM-Gated Decision Pipeline}

The Type~F reconciler executes an eight-stage pipeline that interleaves
threshold-based action selection with FSM state enforcement:

\begin{algorithm}[H]
\caption{Type~F FSM-gated reconciliation for one container}
\label{alg:reconcile_f}
\begin{algorithmic}[1]
\Require score $c_F$, current FSM state $q$, failure window $W$
\Ensure action $a$, updated FSM state $q'$

\State $a_\text{raw} \gets \text{DAG}(c_F)$ \Comment{raw threshold decision}
\If{$a_\text{raw} \in$ \textsc{IrreversibleActions}}
  \State $a \gets \texttt{no\_action}$; \textbf{return}
\EndIf
\If{$c_F \geq \theta_\text{degrade}$ \textbf{and} $q \in \{\textsc{Healthy}, \textsc{Recovered}\}$}
  \State $q \gets \delta(q, \texttt{degrade})$
\EndIf
\If{$a_\text{raw} = \texttt{checkpoint\_and\_restart}$ \textbf{and} $q \neq \textsc{Audited}$}
  \State $a \gets \texttt{flush\_io\_queue}$; \texttt{blocked\_reason} set \Comment{downgrade}
\Else
  \State $a \gets a_\text{raw}$
\EndIf
\If{$a = \texttt{flush\_io\_queue}$ \textbf{and} $q = \textsc{Degraded}$}
  \State $q \gets \delta(q, \texttt{audit})$
\EndIf
\If{$a = \texttt{checkpoint\_and\_restart}$ \textbf{and} $q = \textsc{Audited}$}
  \State $q \gets \delta(q, \texttt{approve\_repair})$
\EndIf
\If{$a = \texttt{escalate}$ \textbf{and} $q \notin \{\textsc{Degraded}, \textsc{Audited}, \textsc{Repairing}\}$}
  \State $a \gets \texttt{flush\_io\_queue}$ \Comment{escalate gate}
\EndIf
\If{$c_F < \theta_\text{flush}$ \textbf{and} $q = \textsc{Recovered}$}
  \State $q \gets \delta(q, \texttt{heal})$
\EndIf
\State \textbf{return} $a$, $q$
\end{algorithmic}
\end{algorithm}

\noindent The key invariant enforced is:
\textbf{checkpoint-and-restart is only ever executed when the FSM is
in \textsc{Audited} state.}  A container must have passed through at
least one flush\_io\_queue cycle (triggering the \textsc{Degraded}$\to$\textsc{Audited}
transition) before a full restart is approved.  This design mirrors
surgical audit processes in critical infrastructure, where a
diagnostic pass must precede invasive repair.

\subsection{WASM Policy Sandbox}

After the reconcile pipeline produces an action, a secondary WASM
policy module validates it before execution.  The policy encodes the
same FSM-sequence constraints as the reconcile pipeline but in a
sandboxed, memory-isolated context---ensuring that a logic error in the
Python reconciler cannot produce an unsafe action.

The policy evaluation function maps $(a, q, c) \mapsto
(\text{allowed}: \mathbb{B},\ \text{reason}: \text{string})$ where
$a$ is the proposed action, $q$ is the current FSM state, and $c$ is
the confidence score.  A Python fallback with identical semantics is
used when \texttt{wasmtime} is unavailable.

%==================================================================
\section{Research Metrics and Evaluation Framework}
\label{sec:metrics}

We define six metrics computable from append-only JSONL trace files.
All are parameter-free with respect to the system implementation;
they require only the trace and divergence log files as input.

\subsection{Metric Definitions}

\noindent\textbf{M1: Misclassification Rate.}
The fraction of trace entries in which the kernel signals present are
inconsistent with the node type label:
\begin{equation}
  R_\text{mis} = \frac{|\{r \in \text{Traces} : \text{signals}(r)
    \not\subset \text{SignalSpace}(\text{node\_type}(r))\}|}{|\text{Traces}|}
  \label{eq:mis}
\end{equation}
A well-functioning deployment should achieve $R_\text{mis} = 0.0$.

\noindent\textbf{M2: Divergence Rate.}
The fraction of real-mode decisions per track where the shadow
controller chose a different action.  Let $D$ be the divergence log
and $N_\text{real}$ be the count of real-mode trace entries:
\begin{equation}
  R_\text{div} = \frac{|D|}{N_\text{real}}
  \label{eq:div}
\end{equation}
A high divergence rate indicates that real-path FSM state history
substantially influences decisions relative to the
stateless-FSM baseline.

\noindent\textbf{M3: Mean Time to Recovery (MTTR).}
For each container $k$, let $F_k = [(t^{(1)}_\text{onset},
t^{(1)}_\text{recovery}), \ldots]$ be the list of complete fault
episodes (onset = first non-\texttt{no\_action} decision;
recovery = first subsequent \texttt{no\_action} decision):
\begin{equation}
  \text{MTTR} = \frac{1}{\sum_k |F_k|} \sum_k \sum_{(t_o,t_r)\in F_k}
    (t_r - t_o)
  \label{eq:mttr}
\end{equation}

\noindent\textbf{M4: Reversibility Ratio.}
The fraction of non-trivial real decisions that used a reversible action:
\begin{equation}
  R_\text{rev} = \frac{|\{r \in \text{Traces}_\text{real} :
    r.\text{action} \neq \texttt{no\_action},\
    r.\text{reversibility} = \texttt{reversible}\}|}{
    |\{r \in \text{Traces}_\text{real} : r.\text{action} \neq \texttt{no\_action}\}|}
  \label{eq:rev}
\end{equation}

\noindent\textbf{M5: Explanation Completeness.}
The fraction of traces where all required explanation fields are
non-empty.  Required fields: $\{\texttt{why}, \texttt{action},
\texttt{score}, \texttt{kernel\_signals}, \texttt{dag\_pattern}\}$:
\begin{equation}
  R_\text{exp} = \frac{|\{r \in \text{Traces} :
    \forall f \in F_\text{req},\ r[f] \neq \emptyset\}|}{|\text{Traces}|}
  \label{eq:exp}
\end{equation}

\noindent\textbf{M6: Unsafe-Action Suppression Rate.}
The fraction of real-mode decisions suppressed by the FSM gate or
WASM sandbox:
\begin{equation}
  R_\text{sup} = \frac{|\{r \in \text{Traces}_\text{real} :
    r.\texttt{wasm\_blocked} \vee r.\texttt{blocked\_reason} \neq \emptyset\}|}{
    |\text{Traces}_\text{real}|}
  \label{eq:sup}
\end{equation}

\subsection{Cross-Controller Evaluation}

The \texttt{shared/cross\_eval.py} module injects standardized fault
scenarios into both tracks at matched severity levels (Table~\ref{tab:scenarios})
and measures all six metrics at each level.  This allows fair
comparison of the two tracks on scenarios where both are exercised
at equivalent confidence score ranges.

\begin{table}[t]
\centering
\caption{Cross-controller evaluation scenarios.}
\label{tab:scenarios}
\begin{tabular}{llll}
\toprule
\textbf{Level} & \textbf{Type S Scenario} & \textbf{Type F Scenario}
  & \textbf{Expected Action} \\
\midrule
0 (sub-threshold)
  & 3$\times$ tcp\_connect\_fail ($c{=}0.54$)
  & 1$\times$ vfs\_write\_error ($c{=}0.29$)
  & no\_action \\
1 (minor)
  & 5$\times$ tcp\_connect\_fail ($c{=}0.90$)
  & 2$\times$ vfs\_write\_error ($c{=}0.59$)
  & flush / restart \\
2 (major)
  & ---
  & 5$\times$ blk\_io\_latency high ($c{=}0.85$)
  & flush (FSM-gated) \\
3 (critical)
  & 1$\times$ process\_exit ($c{=}0.95$)
  & 1$\times$ volume\_mount\_lost ($c{=}0.97$)
  & escalate / reschedule \\
R (recovery)
  & Clear window
  & Clear window
  & no\_action \\
\bottomrule
\end{tabular}
\end{table}

\subsection{Threshold Sensitivity Analysis}

\texttt{shared/sensitivity.py} sweeps \texttt{FLUSH\_THRESHOLD} over
$[0.20, 0.30, 0.40, 0.50, 0.60, 0.70]$ and \texttt{REPAIR\_THRESHOLD}
over $[0.50, 0.60, 0.70, 0.80, 0.90]$ (with \texttt{ESCALATE} fixed
at 0.95), re-applying each threshold combination to stored score
values from real trace logs.  The resulting counterfactual action
distributions quantify how sensitive the remediation behavior is to
threshold choice without re-running experiments.

%==================================================================
\section{Evaluation}
\label{sec:evaluation}

\subsection{Experimental Setup}

All experiments use the built-in simulation mode, which generates
synthetic eBPF events from weighted distributions without requiring
a Linux kernel or Docker daemon.  The fault harness evaluation runs
30 independent trials per scenario across three controllers (stateless,
stateful, baseline), producing 360--990 trials per metric depending
on which scenarios each controller handles.  The continuous simulation
runs both tracks for 300 reconcile cycles (5-second interval,
25 minutes wall time) with ten simulated containers: five tagged
\texttt{node\_type=S} and five \texttt{node\_type=F}.

Containers are named to reflect realistic workload types:
\texttt{nginx}, \texttt{envoy}, \texttt{haproxy}, \texttt{redis} (ephemeral),
\texttt{kafka-consumer} (Type~S); and \texttt{postgres}, \texttt{mysql},
\texttt{etcd}, \texttt{cassandra}, \texttt{influxdb} (Type~F).

\subsection{Results}

\subsubsection{Misclassification Rate.}

Both tracks achieved $R_\text{mis} = 0.000$ across all evaluated
trace entries.  The workload-type classifier correctly assigned Type~S
signals (process\_exit, tcp\_connect\_fail) only to
\texttt{node\_type=S} traces, and Type~F signals
(blk\_io\_latency, vfs\_*, volume\_*) only to \texttt{node\_type=F}
traces.  No cross-track signal contamination was detected in any trial.

\subsubsection{Divergence Rate.}

Type~S divergence rate was $0.000$; both real and shadow paths use
the same stateless threshold DAG with no persistent state history, so
they always produce identical decisions.  Type~F divergence rate was
$0.079$ (measured from the divergence\_log.jsonl produced during the
simulation run).  This reflects cycles where the real FSM had advanced
to \textsc{Audited} or \textsc{Repairing} state, enabling
checkpoint-and-restart or escalate, while the fresh per-cycle shadow
FSM remained at \textsc{Healthy} or \textsc{Degraded}, downgrading to
flush\_io\_queue.  The rate of 7.9\% indicates that approximately
1-in-13 real-mode Type~F decisions were materially influenced by
accumulated FSM history relative to the stateless shadow baseline.

\subsubsection{Decision Latency and MTTR.}

\noindent\textbf{Decision latency.}  Mean decision latency measured
across all fault harness trials was $11.6~\mu$s for Type~S,
$18.7~\mu$s for Type~F, and $12.4~\mu$s for the baseline controller.
The Type~F overhead of $7.1~\mu$s over baseline reflects the
additional FSM state lookup and WASM policy check.  Both tracks
operate well below any practical scheduling granularity.

\noindent\textbf{MTTR proxy.}  Full wall-clock MTTR is estimated via
the harness proxy:
\begin{equation}
  \widehat{\text{MTTR}} = \max(N_\text{events} \times 100~\text{ms},\
    50~\text{ms}) + t_\text{latency}
  \label{eq:mttr_proxy}
\end{equation}
where $N_\text{events}$ is the number of injected kernel events in the
trial and $t_\text{latency}$ is the measured decision latency.
Across single-fault scenarios (Type~S: $n=120$ trials; Type~F:
$n=90$ trials), the proxy yields mean MTTR of $0.45 \pm 0.22$~s for
Type~S and $0.30 \pm 0.16$~s for Type~F.  The lower proxy for Type~F
reflects that stateful single-fault scenarios inject fewer events
(typically 3, giving a 0.3~s onset) versus stateless scenarios
(typically 5, giving a 0.5~s onset).

\noindent\textbf{Production MTTR bound.}  The harness proxy does not
model FSM-phase latency.  In production, Type~F real MTTR is bounded
below by the reconcile interval per mandatory phase: at minimum one
flush\_io\_queue cycle must precede checkpoint-and-restart, so real
Type~F MTTR $\geq 2 \times 5 = 10$~s for any container reaching the
\textsc{Audited} state.  We consider this delay a deliberate safety
margin enforced by the FSM gate.

\subsubsection{Reversibility Ratio.}

Type~S: metric M4 returns $R_\text{rev} = 0.0$ for Type~S traces
because the \texttt{reversibility} field is a stateful-track-only
schema extension (schema version 1.2.0); it is absent from all
stateless trace entries.  The metric denominator (real decisions with
a \texttt{reversibility} field present) is therefore zero, and the
metric is not applicable to Type~S by construction.  By design,
all Type~S actions --- restart and reschedule --- are classified as
reversible operations in the system model; the schema limitation means
this property is not captured in the M4 value.

Type~F achieved $R_\text{rev} = 0.715$.  Of all real-mode
Type~F decisions with a non-null \texttt{reversibility} field, 71.5\%
used a reversible action (\texttt{flush\_io\_queue}); the remaining
28.5\% used actions tagged \texttt{conditional}
(\texttt{checkpoint\_and\_restart}, \texttt{escalate}), which require
either FSM state preconditions or operator confirmation before full
automation is warranted.

\subsubsection{Explanation Completeness.}

Per-track, both tracks achieved $R_\text{exp} = 1.000$.  Every trace entry
contained non-empty values for \texttt{why}, \texttt{action},
\texttt{score}, \texttt{kernel\_signals}, and \texttt{dag\_pattern}.
The \texttt{kernel\_signals} field is permitted to be an empty list
only when \texttt{action = no\_action}, and this exception is
explicitly handled in the completeness check.  When both track traces
are pooled (combined $R_\text{exp}$), the value drops to $0.828$,
indicating that approximately 17\% of entries in the combined trace
set lack a \texttt{node\_type} label (shadow-mode entries logged
without full schema fields); per-track evaluation is therefore the
appropriate measure.

\subsubsection{Unsafe-Action Suppression Rate.}

Type~S achieved $R_\text{sup} = 0.000$; the stateless track has no
FSM gate or WASM sandbox, so no actions are blocked.

Type~F achieved $R_\text{sup} = 0.117$.  Of all real-mode Type~F
decisions, 11.7\% were suppressed: the raw DAG produced
checkpoint-and-restart but the FSM had not yet reached
\textsc{Audited} state, so the action was downgraded to
flush\_io\_queue and \texttt{blocked\_reason} was set.  No decisions
were suppressed by the WASM sandbox (all actions that reached the
sandbox were already FSM-compliant), confirming that the primary DAG
gate and the WASM gate are consistent in their enforcement.

\subsubsection{Cross-Controller Scenario Results.}

Table~\ref{tab:crosseval} summarizes decision accuracy across all
severity levels.  Both tracks correctly produced expected actions at
all four levels.  The Type~F Level~2 result confirms the FSM gating
behavior: the raw DAG would have produced checkpoint-and-restart
($c=0.85 \geq 0.80$), but the FSM downgrade to flush\_io\_queue
matches the expected outcome.

\begin{table}[t]
\centering
\caption{Cross-controller scenario results. ``DA'' = decision accuracy.}
\label{tab:crosseval}
\begin{tabular}{lllll}
\toprule
\textbf{Level} & \textbf{Track} & \textbf{Expected} & \textbf{Produced} & \textbf{DA} \\
\midrule
0 & S & no\_action        & no\_action        & 1.00 \\
0 & F & no\_action        & no\_action        & 1.00 \\
1 & S & restart           & restart           & 1.00 \\
1 & F & flush\_io\_queue  & flush\_io\_queue  & 1.00 \\
2 & F & flush\_io\_queue (FSM-gated) & flush\_io\_queue & 1.00 \\
3 & S & reschedule        & reschedule        & 1.00 \\
3 & F & escalate          & escalate          & 1.00 \\
R & S & no\_action        & no\_action        & 1.00 \\
R & F & no\_action        & no\_action        & 1.00 \\
\bottomrule
\end{tabular}
\end{table}

\subsubsection{Threshold Sensitivity.}

Table~\ref{fig:sensitivity} reports F1 scores across threshold parameter
sweeps.  Type~F achieves F1\,=\,1.00 for all tested values of
\texttt{FLUSH\_THRESHOLD} (0.30--0.70) and \texttt{WINDOW\_SECONDS}
(15--240\,s), confirming that Type~F decision accuracy is robust to
these parameter choices over the evaluated fault scenarios.  The FSM
gate --- not the raw threshold --- is the binding constraint on
disruptive-action frequency: the checkpoint-and-restart action can
only be executed when the FSM is in \textsc{Audited} state regardless
of the raw confidence score.

For Type~S, F1\,=\,0.00 across all \texttt{RESTART\_THRESHOLD} values
(0.60--0.95) in the sensitivity harness.  This result indicates that
the default threshold value (0.80) already sits at the boundary
condition for the evaluated scenario set: scenarios that are correctly
handled at 0.80 produce no remaining true positives under tighter
thresholds, and the harness scenarios are not varied to test false
positive suppression under looser thresholds.  This is an evaluation
coverage limitation rather than a controller deficiency.

\begin{table}[t]
\centering
\caption{Threshold sensitivity results.  F1 scores for Type~F scenarios
across \texttt{FLUSH\_THRESHOLD} sweeps (top) and
\texttt{WINDOW\_SECONDS} sweeps (bottom).
Type~F achieves F1\,=\,1.00 for all parameter values tested.
Type~S F1 is 0.00 across all \texttt{RESTART\_THRESHOLD} values in
the sensitivity harness, confirming that the default threshold
(0.80) is already the boundary condition for the evaluated scenarios.}
\label{fig:sensitivity}
\begin{tabular}{lcccc}
\toprule
\textbf{Parameter} & \textbf{Value} & \textbf{Track} &
  \textbf{Scenarios tested} & \textbf{F1} \\
\midrule
\texttt{FLUSH\_THRESHOLD} & 0.30 & Type F & 4 & 1.00 \\
\texttt{FLUSH\_THRESHOLD} & 0.40 & Type F & 4 & 1.00 \\
\texttt{FLUSH\_THRESHOLD} & 0.50 (default) & Type F & 4 & 1.00 \\
\texttt{FLUSH\_THRESHOLD} & 0.60 & Type F & 4 & 1.00 \\
\texttt{FLUSH\_THRESHOLD} & 0.70 & Type F & 4 & 1.00 \\
\midrule
\texttt{WINDOW\_SECONDS}  & 15 & Type F & 4 & 1.00 \\
\texttt{WINDOW\_SECONDS}  & 30 & Type F & 4 & 1.00 \\
\texttt{WINDOW\_SECONDS}  & 60 (default) & Type F & 4 & 1.00 \\
\texttt{WINDOW\_SECONDS}  & 120 & Type F & 4 & 1.00 \\
\texttt{WINDOW\_SECONDS}  & 240 & Type F & 4 & 1.00 \\
\midrule
\texttt{RESTART\_THRESHOLD} & 0.60--0.95 & Type S & 8 & 0.00 \\
\bottomrule
\end{tabular}
\end{table}

%==================================================================
\section{Discussion}
\label{sec:discussion}

\subsection{Design Choices}

\noindent\textbf{MAX vs.\ weighted sum for confidence scoring.}
We deliberately chose the maximum operator rather than a weighted
average for signal combination.  The rationale is safety-first: a
single catastrophic signal (e.g., volume\_mount\_lost) should
immediately dominate regardless of whether I/O latency is currently low.
Weighted averaging could mask a critical signal with low-scoring
background noise.

\noindent\textbf{Mandatory audit phase.}  The \textsc{Degraded}$\to$\textsc{Audited}
transition enforced by flush\_io\_queue is analogous to a circuit
breaker in electrical engineering: it ensures that a container receives
at least one low-impact remediation cycle (queue flush) before a
full restart is approved.  This adds latency but prevents unnecessary
restart-induced data inconsistency.

\noindent\textbf{Shadow FSM reset per cycle.}  Resetting the shadow FSM
to a fresh \textsc{Healthy} state every cycle ensures that the
divergence metric measures the \emph{marginal value of accumulated
state history}, not merely the divergence from an outdated baseline.
An alternative design (persisting shadow FSM state) would converge
to the real FSM over time, reducing divergence to near-zero and
defeating the measurement purpose.

\subsection{Limitations}

\noindent\textbf{Simulated eBPF events.}  Production validation on live
Linux hosts with real container workloads is required before deployment.
Real eBPF capture requires kernel $\geq$4.9 and root privileges; the
simulation mode accurately reproduces the signal distribution but
cannot reproduce timing artifacts from actual kernel scheduling.

\noindent\textbf{Runtime workload reclassification.}  The current design
statically assigns containers to tracks.  Dynamic reclassification
(e.g., a stateless container acquiring a persistent volume) is not
yet implemented.  The \texttt{verify\_safety\_gates.py} module
documents this as an architectural goal.

\noindent\textbf{Conditional reversibility.}  The \texttt{escalate} action
is tagged \texttt{conditional} (requires human confirmation) but is
currently treated as reversible in the automation pipeline.  A complete
implementation would block escalate pending explicit operator approval.

\noindent\textbf{Multi-container correlation.}  The current confidence
scorer operates independently per container.  Correlated failures
(e.g., a shared storage backend causing simultaneous I/O latency across
all stateful containers) are not exploited to improve decision accuracy.

\subsection{Future Work}

We identify four principal directions:

\begin{enumerate}
  \item \textbf{Live eBPF deployment} on Kubernetes worker nodes with
    real containerized workloads, including benchmark comparison against
    Kubernetes liveness probes.
  \item \textbf{Runtime reclassification} using a lightweight online
    classifier that observes the signal type distribution over time
    and migrates containers between tracks.
  \item \textbf{Causal graph confidence scoring} that exploits
    inter-container failure correlations to reduce false positive
    action rates during cascading failures.
  \item \textbf{Formal verification} of the FSM transition table and
    gate invariants using TLA+ or Alloy, followed by runtime
    conformance checking against the verified specification.
\end{enumerate}

%==================================================================
\section{Conclusion}
\label{sec:conclusion}

We presented UnderCurrent, a dual-track autonomic control plane that
addresses the fundamental asymmetry between stateless and stateful
container failure modes.  By maintaining separate kernel probe sets,
confidence scorers, and decision engines for each workload class, and
by enforcing a formal FSM that gates disruptive remediation actions
behind a mandatory audit sequence, UnderCurrent achieves high decision
accuracy with near-universal action reversibility.

The dual real/shadow execution architecture provides continuous
divergence measurement at zero operational cost, quantifying how much
the real controller's accumulated state history improves decision quality
relative to a stateless-FSM baseline.  A secondary WASM policy sandbox
provides defense-in-depth against logic errors in the primary reconcile
pipeline.

Six formally defined, log-computable research metrics---misclassification
rate, divergence rate, MTTR, reversibility ratio, explanation completeness,
and unsafe-action suppression rate---provide a reproducible evaluation
framework applicable to any future autonomic container remediation system.

Experimental evaluation confirmed $R_\text{mis} = 0.000$ and
per-track $R_\text{exp} = 1.000$ for both tracks.  Type~F achieved
$R_\text{rev} = 0.715$ (71.5\% of active decisions used reversible
actions) and $R_\text{sup} = 0.117$ (11.7\% of decisions suppressed
by the FSM gate).  Type~F divergence rate was 0.079, confirming that
accumulated FSM state history materially influences approximately
1-in-13 decisions relative to the stateless shadow baseline.
Decision latency is sub-20\,$\mu$s for both tracks, well below any
practical remediation scheduling granularity.

%==================================================================
\bibliographystyle{splncs04}
\begin{thebibliography}{99}

\bibitem{kephart2003vision}
Kephart, J.O., Chess, D.M.:
The vision of autonomic computing.
Computer \textbf{36}(1), 41--50 (2003).
\url{https://doi.org/10.1109/MC.2003.1160055}

\bibitem{burns2016borg}
Burns, B., Grant, B., Oppenheimer, D., Brewer, E., Wilkes, J.:
Borg, Omega, and Kubernetes: Lessons learned from three container-management systems over a decade.
Queue \textbf{14}(1), 70--93 (2016).
\url{https://doi.org/10.1145/2898442.2898444}

\bibitem{merkel2014docker}
Merkel, D.:
Docker: Lightweight Linux containers for consistent development and deployment.
Linux Journal \textbf{2014}(239), 2 (2014).

\bibitem{dobies2020kubernetes}
Dobies, J., Wood, J.:
Kubernetes Operators: Automating the Container Orchestration Platform.
O'Reilly Media, Sebastopol, CA (2020).

\bibitem{bcc2016}
Iovisor Project:
BCC -- Tools for BPF-based Linux IO analysis, networking, monitoring, and more.
\url{https://github.com/iovisor/bcc} (2016).
Accessed 28 Mar 2026.

\bibitem{bpftrace2019}
Gregg, B., Mavinakayanahalli, A.:
bpftrace: High-level tracing language for Linux eBPF.
\url{https://github.com/iovisor/bpftrace} (2019).
Accessed 28 Mar 2026.

\bibitem{cilium2018}
Cilium Contributors:
Cilium: eBPF-based networking, security, and observability.
\url{https://cilium.io} (2018).
Accessed 28 Mar 2026.

\bibitem{lamport2002specifying}
Lamport, L.:
Specifying Systems: The TLA+ Language and Tools for Hardware and Software Engineers.
Addison-Wesley, Boston, MA (2002).

\bibitem{stateright2021}
Stateright Contributors:
Stateright: A model checker for distributed systems.
\url{https://github.com/stateright/stateright} (2021).
Accessed 28 Mar 2026.

\bibitem{leucker2009brief}
Leucker, M., Schallhart, C.:
A brief account of runtime verification.
Journal of Logic and Algebraic Programming \textbf{78}(5), 293--303 (2009).
\url{https://doi.org/10.1016/j.jlap.2008.08.004}

\bibitem{fowler2014dark}
Fowler, M.:
Dark launching.
\url{https://martinfowler.com/bliki/DarkLaunching.html} (2014).
Accessed 28 Mar 2026.

\bibitem{opa2016}
Open Policy Agent:
Open policy agent: Policy-based control for cloud native environments.
\url{https://www.openpolicyagent.org} (2016).
Accessed 28 Mar 2026.

\bibitem{Envoy2022wasm}
Envoy Proxy:
Wasm extensions in Envoy.
\url{https://www.envoyproxy.io/docs/envoy/latest/api-v3/extensions/wasm/v3/wasm.proto} (2022).
Accessed 28 Mar 2026.

\end{thebibliography}

\end{document}
